\chapter{2003年北京卷（秋理）}
\section{单选题}
\begin{problem}
设集合 \(A=\{x\mid x^2-1>0\}\)，\(B=\{x\mid \log_2 x>0\}\)，\(A\cap B\) 等于（\quad）

\choices
    {\(\{x\mid x>1\}\)}
    {\(\{x\mid x>0\}\)}
    {\(\{x\mid x<-1\}\)}
    {\(\{x\mid x<-1\text{\ 或\ }x>1\}\)}
\answer{A}

\begin{solution}
由 \(x^2-1>0\) 得 \(x<-1\) 或 \(x>1\)，所以 \(A=(-\infty,-1)\cup(1,+\infty)\). 又因为底数 \(2>1\)，\(\log_2x>0\) 等价于 \(x>1\)，故 \(B=(1,+\infty)\). 因此 \(A\cap B=(1,+\infty)=\{x\mid x>1\}\)，应选 A.
\end{solution}
\end{problem}

\begin{problem}
设 \(y_1=4^{0.9}\)，\(y_2=8^{0.44}\)，\(y_3=\left(\frac{1}{2}\right)^{-1.5}\)，则（\quad）

\choices
    {\(y_3>y_1>y_2\)}
    {\(y_2>y_1>y_3\)}
    {\(y_1>y_2>y_3\)}
    {\(y_1>y_3>y_2\)}
\answer{D}

\begin{solution}
把各数化为底数为 \(2\) 的幂，有 \(y_1=4^{0.9}=2^{1.8}\)，\(y_2=8^{0.44}=2^{1.32}\)，\(y_3=\left(\frac12\right)^{-1.5}=2^{1.5}\). 由于 \(2>1\)，且 \(1.8>1.5>1.32\)，所以 \(y_1>y_3>y_2\)，应选 D.
\end{solution}
\end{problem}

\begin{problem}
“\(\cos 2\alpha=-\frac{\sqrt{3}}{2}\)”是“\(\alpha=k\pi+\frac{5\pi}{12}\)，\(k\in\Z\)”的（\quad）

\choices
    {必要非充分条件}
    {充分非必要条件}
    {充分必要条件}
    {既非充分又非必要条件}
\answer{A}

\begin{solution}
由 \(\cos 2\alpha=-\frac{\sqrt3}{2}\)，得 \(2\alpha=\frac{5\pi}{6}+2k\pi\) 或 \(2\alpha=\frac{7\pi}{6}+2k\pi\)，其中 \(k\in\Z\). 因此 \(\alpha=\frac{5\pi}{12}+k\pi\) 或 \(\alpha=\frac{7\pi}{12}+k\pi\). 后一个条件“\(\alpha=k\pi+\frac{5\pi}{12}\)”能推出前一个等式，但前一个等式还包含 \(\alpha=k\pi+\frac{7\pi}{12}\) 的情形，所以它是必要非充分条件，应选 A.
\end{solution}
\end{problem}

\begin{problem}
已知 \(\alpha\)，\(\beta\) 是平面，\(m\)，\(n\) 是直线. 下列命题中不正确的是（\quad）

\choices
    {若 \(m\parallel n\)，\(m\perp\alpha\)，则 \(n\perp\alpha\)}
    {若 \(m\parallel\alpha\)，\(\alpha\cap\beta=n\)，则 \(m\parallel n\)}
    {若 \(m\perp\alpha\)，\(m\perp\beta\)，则 \(\alpha\parallel\beta\)}
    {若 \(m\perp\alpha\)，\(m\subset\beta\)，则 \(\alpha\perp\beta\)}
\answer{B}

\begin{solution}
若 \(m\parallel n\) 且 \(m\perp\alpha\)，则与 \(m\) 平行的直线 \(n\) 也垂直于平面 \(\alpha\)，所以命题 A 正确. 若两平面都垂直于同一直线 \(m\)，则它们的法向方向相同，从而 \(\alpha\parallel\beta\)，命题 C 正确. 若 \(m\perp\alpha\) 且 \(m\subset\beta\)，平面 \(\beta\) 内有一条直线垂直于平面 \(\alpha\)，所以 \(\alpha\perp\beta\)，命题 D 正确.

命题 B 不成立：例如取 \(\alpha\) 为平面 \(z=0\)，\(\beta\) 为平面 \(y=0\)，则 \(n=\alpha\cap\beta\) 为 \(x\) 轴；再取直线 \(m\colon y=x\)，\(\ z=1\)，则 \(m\parallel\alpha\)，但 \(m\not\parallel n\). 因此不正确的是 B.
\end{solution}
\end{problem}

\begin{problem}
极坐标方程 \(\rho^2\cos 2\theta-2\rho\cos\theta=1\) 表示的曲线是（\quad）

\choices
    {圆}
    {椭圆}
    {抛物线}
    {双曲线}
\answer{D}

\begin{solution}
令 \(x=\rho\cos\theta\)，\(y=\rho\sin\theta\)，则 \(\rho^2\cos2\theta=x^2-y^2\)，且 \(\rho\cos\theta=x\). 原方程化为 \(x^2-y^2-2x=1\)，即 \((x-1)^2-y^2=2\). 这是以 \((1,0)\) 为中心、实轴平行于 \(x\) 轴的双曲线，应选 D.
\end{solution}
\end{problem}

\begin{problem}
若 \(z\in\mathbb{C}\) 且 \(|z+2-2\i|=1\)，则 \(|z-2-2\i|\) 的最小值是（\quad）

\choices
    {\(2\)}
    {\(3\)}
    {\(4\)}
    {\(5\)}
\answer{B}

\begin{solution}
令 \(w=z+2-2\i\)，则 \(\lvert w\rvert=1\)，并且 \(z-2-2\i=w-4\). 由三角不等式的逆不等式，\(\lvert z-2-2\i\rvert=\lvert w-4\rvert\geqslant\bigl|\lvert4\rvert-\lvert w\rvert\bigr|=3\). 取 \(w=1\) 时等号成立，所以所求最小值为 \(3\)，应选 B.
\end{solution}
\end{problem}

\begin{problem}
如果圆台的母线与底面成 \(60^{\circ}\) 角，那么这个圆台的侧面积与轴截面面积的比为（\quad）

\choices
    {\(2\pi\)}
    {\(\frac{3}{2}\pi\)}
    {\(\frac{2\sqrt{3}}{3}\pi\)}
    {\(\frac{1}{2}\pi\)}
\answer{C}

\begin{solution}
设圆台上、下底面的半径分别为 \(R\)、\(r_0\)，母线长为 \(\ell\)，高为 \(h\). 母线与底面成 \(60^{\circ}\) 角，所以 \(h=\ell\sin60^{\circ}\). 圆台的侧面积为 \(S_1=\pi(R+r_0)\ell\)，轴截面是高为 \(h\)、两底边长分别为 \(2R\)、\(2r_0\) 的等腰梯形，故其面积为 \(S_2=(R+r_0)h\). 因此
\[
\frac{S_1}{S_2}=\frac{\pi\ell}{h}=\frac{\pi}{\sin60^{\circ}}=\frac{2\sqrt3}{3}\pi
\]
所以应选 C.
\end{solution}
\end{problem}

\begin{problem}
从黄瓜，白菜，油菜，扁豆 \(4\text{ 种}\)蔬菜品种中选出 \(3\text{ 种}\)，分别种在不同土质的三块土地上，其中黄瓜必须种植，不同的种植方法共有（\quad）

\choices
    {\(24\text{ 种}\)}
    {\(18\text{ 种}\)}
    {\(12\text{ 种}\)}
    {\(6\text{ 种}\)}
\answer{B}

\begin{solution}
黄瓜必须选取，再从白菜、油菜、扁豆中选取 \(2\) 种，有 \(\mathrm C_3^2=3\) 种选法. 选定 \(3\) 种后，将它们分别种在三块不同土质的土地上，有 \(\mathrm A_3^3=3!=6\) 种安排. 因此不同种植方法共有 \(3\times6=18\) 种，应选 B.
\end{solution}
\end{problem}

\begin{problem}
若数列 \(\{a_n\}\) 的通项公式是
\(a_n=\frac{3^{-n}+2^{-n}+(-1)^n\left(3^{-n}-2^{-n}\right)}{2}\)（\(n=1\)，\(2\)，\(\dots\)），则 \(\displaystyle\lim_{n\to\infty}(a_1+a_2+\cdots+a_n)\) 等于（\quad）

\choices
    {\(\frac{11}{24}\)}
    {\(\frac{17}{24}\)}
    {\(\frac{19}{24}\)}
    {\(\frac{25}{24}\)}
\answer{C}

\begin{solution}
当 \(n\) 为奇数时，\((-1)^n=-1\)，所以 \(a_n=2^{-n}\)；当 \(n\) 为偶数时，\((-1)^n=1\)，所以 \(a_n=3^{-n}\). 因此
\[
\lim_{n\to\infty}\sum_{j=1}^{n}a_j
=\frac{\frac12}{1-\frac14}+\frac{\frac19}{1-\frac19}
=\frac23+\frac18=\frac{19}{24}
\]
故应选 C.
\end{solution}
\end{problem}

\begin{problem}
某班试用电子投票系统选举班干部候选人. 全班 \(k\text{ 名}\)同学都有选举权和被选举权，他们的编号分别为 \(1\)，\(2\)，\(\dots\)，\(k\). 规定：同意按“1”，不同意（含弃权）按“0”，令
\(a_{ij}=\begin{cases}
1, & \text{第\ }i\text{\ 号同学同意第 }j\text{\ 号同学当选},\\
0, & \text{第\ }i\text{\ 号同学不同意第 }j\text{\ 号同学当选},
\end{cases}\)
其中 \(i=1\)，\(2\)，\(\dots\)，\(k\)，且 \(j=1\)，\(2\)，\(\dots\)，\(k\)，则同时同意第 \(1\)，\(2\) 号同学当选的人数为（\quad）

\choices
    {\(a_{11}+a_{12}+\cdots+a_{1k}+a_{21}+a_{22}+\cdots+a_{2k}\)}
    {\(a_{11}+a_{21}+\cdots+a_{1k}+a_{12}+a_{22}+\cdots+a_{k2}\)}
    {\(a_{11}a_{12}+a_{21}a_{22}+\cdots+a_{k1}a_{k2}\)}
    {\(a_{11}a_{21}+a_{12}a_{22}+\cdots+a_{1k}a_{2k}\)}
\answer{C}

\begin{solution}
对固定的第 \(i\) 号同学，\(a_{i1}a_{i2}=1\) 当且仅当他同时同意第 \(1\)、\(2\) 号同学当选，否则该乘积为 \(0\). 所以把所有同学的这一指标相加，所得人数为
\[
\sum_{i=1}^{k}a_{i1}a_{i2}=a_{11}a_{12}+a_{21}a_{22}+\cdots+a_{k1}a_{k2}
\]
因此应选 C.
\end{solution}
\end{problem}

\section{填空题}
\begin{problem}
函数 \(f(x)=\lg(1+x^2)\)，\(g(x)=
\begin{cases}
x+2, & x<-1,\\
0, & |x|\leqslant 1,\\
-x+2, & x>1,
\end{cases}\)
\(h(x)=\tan 2x\) 中，\fillinblank{} 是偶函数.
\answer{\(f(x)\)，\(g(x)\)}

\begin{solution}
对任意定义域内的 \(x\)，有 \(f(-x)=\lg(1+(-x)^2)=\lg(1+x^2)=f(x)\)，所以 \(f\) 是偶函数. 对 \(g\) 分段考察：当 \(x>1\) 时，\(g(x)=-x+2\)，且 \(-x<-1\)，故 \(g(-x)=(-x)+2=-x+2=g(x)\)；当 \(x<-1\) 时，\(-x>1\)，故 \(g(-x)=-(-x)+2=x+2=g(x)\)；当 \(\lvert x\rvert\leqslant1\) 时，\(\lvert -x\rvert\leqslant1\)，从而 \(g(-x)=g(x)=0\). 所以 \(g\) 也是偶函数.

函数 \(h\) 的定义域为 \(\{x\mid \cos2x\ne0\}\)，关于原点对称；在该定义域内，\(h(-x)=\tan(-2x)=-\tan2x=-h(x)\)，所以 \(h\) 是奇函数而不是偶函数. 因此填 \(f(x)\)、\(g(x)\).
\end{solution}
\end{problem}

\begin{problem}
以双曲线 \(\frac{x^2}{16}-\frac{y^2}{9}=1\) 右顶点为顶点，左焦点为焦点的抛物线的方程是\fillinblank{}.
\answer{\(y^2=-36(x-4)\)}

\begin{solution}
双曲线的半长轴为 \(4\)，半短轴为 \(3\)，所以半焦距为 \(c=\sqrt{4^2+3^2}=5\). 因而双曲线的右顶点为 \(V(4,0)\)，左焦点为 \(F(-5,0)\). 设所求抛物线的标准方程为 \(y^2=4p(x-4)\)，则其焦点为 \((4+p,0)\). 由 \(4+p=-5\) 得 \(p=-9\)，所以所求抛物线的方程为 \(y^2=-36(x-4)\).
\end{solution}
\end{problem}

\begin{problem}
如图，已知底面半径为 \(r\) 的圆柱被一个平面所截，剩下部分母线长的最大值为 \(a\)，最小值为 \(b\)，那么圆柱被截后剩下部分的体积是\fillinblank{}.

\bitmapfigure[max width=6.0cm]{img/2003/beijing_science_autumn/q13_fig1.png}
\answer{\(\frac{1}{2}\pi r^{2}(a+b)\)}

\begin{solution}
取圆柱底面为 \(D=\{(x,y)\mid x^2+y^2\leqslant r^2\}\)，以底面为 \(z=0\). 截平面下方剩余部分在底面点 \((x,y)\) 处的母线长度可写成一次函数 \(\ell(x,y)=ux+vy+w\). 若 \((u,v)\ne(0,0)\)，则一次函数在圆盘 \(D\) 上的最大值和最小值分别在一对相反点 \((x_0,y_0)\)、\((-x_0,-y_0)\) 处取得，因此 \(a=\ell(x_0,y_0)\)，\(b=\ell(-x_0,-y_0)\)，从而 \(a+b=2w\)；若 \(u=v=0\)，则 \(\ell\) 为常数，结论同样成立.

所以剩余部分的体积为
\[
V=\iint_D\ell(x,y)\,\mathrm d x\,\mathrm d y
=u\iint_Dx\,\mathrm d x\,\mathrm d y
+v\iint_Dy\,\mathrm d x\,\mathrm d y
+w\iint_D\mathrm d x\,\mathrm d y
\]
前两个积分因圆盘关于原点对称而为 \(0\)，最后一个积分为 \(w\pi r^2\)，故 \(V=w\pi r^2=\frac12\pi r^2(a+b)\).
\end{solution}
\end{problem}

\begin{problem}
将长度为 1 的铁丝分成两段，分别围成一个正方形和一个圆形，要使正方形与圆的面积之和最小，正方形的周长应为\fillinblank{}.
\answer{\(\frac{4}{4+\pi}\)}

\begin{solution}
设正方形的周长为 \(x\)，则圆的周长为 \(1-x\)，其中 \(0<x<1\). 正方形和圆的面积之和为
\[
S(x)=\frac{x^2}{16}+\frac{(1-x)^2}{4\pi}
\]
求导得 \(S'(x)=\frac{x}{8}-\frac{1-x}{2\pi}\)，令 \(S'(x)=0\)，得到 \(\pi x=4(1-x)\)，即 \(x=\frac{4}{4+\pi}\). 又 \(S''(x)=\frac18+\frac{1}{2\pi}>0\)，所以此时面积之和取得最小值. 因此正方形的周长为 \(\frac{4}{4+\pi}\).
\end{solution}
\end{problem}

\section{解答题}
\begin{problem}
已知函数 \(f(x)=\cos^{4}x-2\sin x\cos x-\sin^{4}x\).
\begin{enumerate}
\item 求 \(f(x)\) 的最小正周期；
\item 若 \(x\in\left[0,\frac{\pi}{2}\right]\)，求 \(f(x)\) 的最大值，最小值.
\end{enumerate}

\begin{solution}
\begin{enumerate}
\item 由 \(\cos^4x-\sin^4x=(\cos^2x-\sin^2x)(\cos^2x+\sin^2x)=\cos 2x\)，且 \(2\sin x\cos x=\sin 2x\)，得 \(f(x)=\cos 2x-\sin 2x=\sqrt{2}\cos\left(2x+\frac{\pi}{4}\right)\). 因为自变量 \(x\) 的系数为 \(2\)，所以 \(f(x)\) 的最小正周期为 \(T=\frac{2\pi}{2}=\pi\).
\item 令 \(t=2x\)，则 \(t\in[0,\pi]\)，并设 \(\varphi(t)=\cos t-\sin t\)，于是 \(f(x)=\varphi(t)\). 有 \(\varphi'(t)=-\sin t-\cos t\)，在区间 \([0,\pi]\) 内令 \(\varphi'(t)=0\)，得 \(t=\frac{3\pi}{4}\). 比较端点和驻点的函数值：\(\varphi(0)=1\)，\(\varphi(\frac{3\pi}{4})=-\sqrt{2}\)，\(\varphi(\pi)=-1\). 因而当 \(t=0\)，即 \(x=0\) 时，\(f(x)\) 取得最大值 \(1\)；当 \(t=\frac{3\pi}{4}\)，即 \(x=\frac{3\pi}{8}\) 时，\(f(x)\) 取得最小值 \(-\sqrt{2}\).
\end{enumerate}
\end{solution}
\end{problem}

\begin{problem}
已知数列 \(\{a_n\}\) 是等差数列，且 \(a_1=2\)，\(a_1+a_2+a_3=12\).
\begin{enumerate}
\item 求数列 \(\{a_n\}\) 的通项公式；
\item 令 \(b_n=a_nx^n\)（\(x\in\R\)）. 求数列 \(\{b_n\}\) 前 \(n\) 项和的公式.
\end{enumerate}

\begin{solution}
\begin{enumerate}
\item 设等差数列的公差为 \(d\)，则 \(a_2=2+d\)，\(a_3=2+2d\). 由 \(a_1+a_2+a_3=12\) 得 \(2+(2+d)+(2+2d)=12\)，即 \(6+3d=12\)，所以 \(d=2\). 因此 \(a_n=a_1+(n-1)d=2+2(n-1)=2n\).
\item 由上一问 \(b_n=2nx^n\). 设前 \(n\) 项和为 \(S_n\)，则 \(S_n=2x+4x^2+\cdots+2nx^n\). 当 \(x\ne1\) 时，将等式两边乘以 \(x\)，得 \(xS_n=2x^2+4x^3+\cdots+2nx^{n+1}\)，两式相减可得 \((1-x)S_n=2(x+x^2+\cdots+x^n-nx^{n+1})\). 又 \(x+x^2+\cdots+x^n=\frac{x(1-x^n)}{1-x}\)，所以 \(S_n=\frac{2\left[x-(n+1)x^{n+1}+nx^{n+2}\right]}{(1-x)^2}\). 当 \(x=1\) 时，\(S_n=2(1+2+\cdots+n)=n(n+1)\). 因而前 \(n\) 项和为
\[
S_n=\begin{cases}
\dfrac{2\left[x-(n+1)x^{n+1}+nx^{n+2}\right]}{(1-x)^2}, & x\ne1,\\
n(n+1), & x=1.
\end{cases}
\]
\end{enumerate}
\end{solution}
\end{problem}

\begin{problem}
如图，正三棱柱 \(ABC-A_1B_1C_1\) 的底面边长为 \(3\)，侧棱 \(AA_1=\frac{3\sqrt{3}}{2}\)，\(D\) 是 \(CB\) 延长线上一点，且 \(BD=BC\).

\begin{enumerate}
\item 求证：直线 \(BC_1\parallel\) 平面 \(AB_1D\)；
\item 求二面角 \(B_1-AD-B\) 的大小；
\item 求三棱锥 \(C_1-ABB_1\) 的体积.
\end{enumerate}

\bitmapfigure[max width=4.8cm]{img/2003/beijing_science_autumn/q17_fig1.png}
\begin{solution}
令 \(h=\frac{3\sqrt{3}}{2}\)，建立直角坐标系，使 \(B(0,0,0)\)，\(C(3,0,0)\)，\(A(\frac32,h,0)\)，且棱柱沿 \(z\) 轴方向延伸. 由 \(C\)，\(B\)，\(D\) 依次在同一直线上且 \(BD=BC=3\)，得 \(D(-3,0,0)\)，又 \(B_1(0,0,h)\)，\(C_1(3,0,h)\).
\begin{enumerate}
\item 由坐标直接得到 \(\overrightarrow{BC_1}=(3,0,h)=\overrightarrow{DB_1}\). 直线 \(DB_1\) 在平面 \(AB_1D\) 内，且平面 \(AB_1D\) 的方程为 \(h(x+3)-\frac92y-3z=0\)，点 \(B\) 不在此平面内. 因此 \(BC_1\parallel DB_1\)，且 \(BC_1\) 平行于平面 \(AB_1D\).
\item 平面 \(ADB\) 是底面所在的平面，其一个法向量为 \(\bs n_2=(0,0,1)\). 平面 \(ADB_1\) 内取 \(\overrightarrow{DA}=(\frac92,h,0)\)，\(\overrightarrow{DB_1}=(3,0,h)\)，于是其一个法向量为 \(\bs n_1=\overrightarrow{DA}\times\overrightarrow{DB_1}=(h^2,-\frac92h,-3h)\). 代入 \(h=\frac{3\sqrt3}{2}\) 得 \(\lVert\bs n_1\rVert=9\sqrt3\)，且 \(\lvert\bs n_1\cdot\bs n_2\rvert=3h=\frac{9\sqrt3}{2}\). 设二面角 \(B_1-AD-B\) 的平面角为 \(\theta\)，则 \(\cos\theta=\frac{\lvert\bs n_1\cdot\bs n_2\rvert}{\lVert\bs n_1\rVert\lVert\bs n_2\rVert}=\frac12\)，所以 \(\theta=60^{\circ}\).
\item 三角形 \(ABB_1\) 是直角三角形，故 \(S_{\triangle ABB_1}=\frac12\cdot AB\cdot BB_1=\frac12\cdot3\cdot h=\frac{9\sqrt3}{4}\). 平面 \(ABB_1\) 是过 \(AB\) 的竖直平面，点 \(C_1\) 到该平面的距离等于底面内点 \(C\) 到边 \(AB\) 的距离，即正三角形 \(ABC\) 的高 \(h\). 因此三棱锥 \(C_1-ABB_1\) 的体积为 \(V=\frac13S_{\triangle ABB_1}h=\frac13\cdot\frac{9\sqrt3}{4}\cdot\frac{3\sqrt3}{2}=\frac{27}{8}\).
\end{enumerate}
\end{solution}
\end{problem}

\begin{problem}
如图，椭圆的长轴 \(A_1A_2\) 与 \(x\) 轴平行，短轴 \(B_1B_2\) 在 \(y\) 轴上，中心为 \(M(0,r)\)（\(b>r>0\)）.

\begin{enumerate}
\item 写出椭圆的方程，求椭圆的焦点坐标及离心率；
\item 直线 \(y=k_1x\) 交椭圆于两点 \(C(x_1,y_1)\)，\(D(x_2,y_2)\)（\(y_2>0\)）；直线 \(y=k_2x\) 交椭圆于两点 \(G(x_3,y_3)\)，\(H(x_4,y_4)\)（\(y_4>0\)）. 求证：\(\frac{k_1x_1x_2}{x_1+x_2}=\frac{k_2x_3x_4}{x_3+x_4}\)；
\item 对于前一问中的 \(C\)，\(D\)，\(G\)，\(H\)，设 \(CH\) 交 \(x\) 轴于点 \(P\)，\(GD\) 交 \(x\) 轴于点 \(Q\). 求证：\(|OP|=|OQ|\).（证明过程不考虑 \(CH\) 或 \(GD\) 垂直于 \(x\) 轴的情形）
\end{enumerate}

\bitmapfigure[max width=4.5cm]{img/2003/beijing_science_autumn/q18_fig1.png}
\begin{solution}
\begin{enumerate}
\item 由长轴、短轴的方向知 \(a>b\)，结合题设 \(b>r>0\)，故 \(a>b>r>0\). 椭圆中心为 \(M(0,r)\)，半长轴为 \(a\)，半短轴为 \(b\)，所以方程为 \(\displaystyle\frac{x^2}{a^2}+\frac{(y-r)^2}{b^2}=1\). 设 \(c=\sqrt{a^2-b^2}\)，则两个焦点为 \(F_1(-c,r)\)、\(F_2(c,r)\)，离心率为 \(e=\frac ca=\frac{\sqrt{a^2-b^2}}a\).
\item 对任意斜率为 \(k\) 的直线 \(y=kx\)，代入椭圆方程并整理，得 \((b^2+a^2k^2)x^2-2a^2rkx+a^2(r^2-b^2)=0\). 由于 \(y_2=k_1x_2>0\)、\(y_4=k_2x_4>0\)，有 \(k_1\ne0\)、\(k_2\ne0\)，下列分式均有意义. 当 \(k=k_1\) 时，由韦达定理有 \(x_1+x_2=\frac{2a^2rk_1}{b^2+a^2k_1^2}\)，\(x_1x_2=\frac{a^2(r^2-b^2)}{b^2+a^2k_1^2}\)，从而 \(\displaystyle\frac{k_1x_1x_2}{x_1+x_2}=\frac{r^2-b^2}{2r}\). 当 \(k=k_2\) 时，同一方程给出 \(x_3+x_4=\frac{2a^2rk_2}{b^2+a^2k_2^2}\)，\(x_3x_4=\frac{a^2(r^2-b^2)}{b^2+a^2k_2^2}\)，从而 \(\displaystyle\frac{k_2x_3x_4}{x_3+x_4}=\frac{r^2-b^2}{2r}\). 因此 \(\displaystyle\frac{k_1x_1x_2}{x_1+x_2}=\frac{k_2x_3x_4}{x_3+x_4}\).
\item 设 \(P=(p,0)\)，\(Q=(q,0)\). 由两点式求直线 \(CH\) 与 \(x\) 轴交点的横坐标，得 \(\displaystyle p=\frac{x_1y_4-x_4y_1}{y_4-y_1}=\frac{x_1x_4(k_2-k_1)}{k_2x_4-k_1x_1}\). 由于 \(D\)，\(Q\)，\(G\) 三点共线，有 \(\displaystyle\frac{x_2-q}{x_3-q}=\frac{y_2}{y_3}=\frac{k_1x_2}{k_2x_3}\)，解得 \(\displaystyle q=\frac{x_3y_2-x_2y_3}{y_2-y_3}=\frac{x_2x_3(k_1-k_2)}{k_1x_2-k_2x_3}\). 由上一问的等式可得 \(k_1x_1x_2(x_3+x_4)=k_2x_3x_4(x_1+x_2)\)，整理为 \(\displaystyle\frac{x_1x_4}{k_2x_4-k_1x_1}=\frac{x_2x_3}{k_1x_2-k_2x_3}\). 结合 \(p\) 与 \(q\) 的表达式，得到 \(p=-q\)，所以 \(\lvert OP\rvert=\lvert OQ\rvert\).
\end{enumerate}
\end{solution}
\end{problem}

\begin{problem}
有三个新兴城镇，分别位于 \(A\)，\(B\)，\(C\) 三点处，且 \(AB=AC=a\)，\(BC=2b\). 今计划合建一个中心医院，为同时方便三镇，准备建在 \(BC\) 的垂直平分线上的 \(P\) 点处.（建立坐标系如图）

\begin{enumerate}
\item 若希望点 \(P\) 到三镇距离的平方和为最小，点 \(P\) 应位于何处？
\item 若希望点 \(P\) 到三镇的最远距离为最小，点 \(P\) 应位于何处？
\end{enumerate}

\bitmapfigure[max width=6.0cm]{img/2003/beijing_science_autumn/q19_fig1.png}
\begin{solution}
由三角形的非退化性有 \(a>b>0\). 设 \(h=\sqrt{a^2-b^2}>0\)，则由题设可取 \(A(0,h)\)，\(B(-b,0)\)，\(C(b,0)\)，并设 \(P(0,t)\).
\begin{enumerate}
\item 点 \(P\) 到三镇距离的平方和为 \(L(t)=(t-h)^2+(t^2+b^2)+(t^2+b^2)=3t^2-2ht+h^2+2b^2=3\left(t-\frac h3\right)^2+\frac{2h^2}{3}+2b^2\). 因此当 \(t=\frac h3\) 时平方和最小，即 \(P\) 位于 \(AO\) 上且 \(OP=\frac h3=\frac{\sqrt{a^2-b^2}}3\).
\item 记点 \(P\) 到三镇的最远距离为 \(F(t)=\max\left\{|t-h|,\sqrt{t^2+b^2}\right\}\). 当 \(t<0\) 时，\(|t-h|>h\) 且 \(\sqrt{t^2+b^2}>b\)，故不可能优于 \(t=0\)；当 \(t>h\) 时，两项均随 \(t\) 增大而增大，故最优点在 \([0,h]\) 内. 在 \([0,h]\) 上，\(h-t\) 单调递减，\(\sqrt{t^2+b^2}\) 单调递增，二者相等的条件为 \((h-t)^2=t^2+b^2\)，即 \(t=\frac{h^2-b^2}{2h}\). 若 \(h\leq b\)，该值不为正，且在 \([0,h]\) 上始终有 \(h-t\leq\sqrt{t^2+b^2}\)，所以 \(F(t)\) 在 \(t=0\) 时最小，\(P=O\). 若 \(h>b\)，令 \(t_0=\frac{h^2-b^2}{2h}\)，则 \(0<t_0<h\)，在 \(0\leq t\leq t_0\) 时 \(F(t)=h-t\) 递减，在 \(t_0\leq t\leq h\) 时 \(F(t)=\sqrt{t^2+b^2}\) 递增，故最小点为 \(P(0,t_0)\). 由于 \(h^2=a^2-b^2\)，这等价于：当 \(a^2\leq2b^2\) 时 \(P=O\)；当 \(a^2>2b^2\) 时 \(P\) 位于 \(AO\) 上，且 \(\displaystyle OP=\frac{a^2-2b^2}{2\sqrt{a^2-b^2}}\).
\end{enumerate}
\end{solution}
\end{problem}

\begin{problem}
设 \(y=f(x)\) 是定义在区间 \([-1,1]\) 上的函数，且满足条件：
\begin{circlelist}
\item \(f(-1)=f(1)=0\)；
\item 对任意的 \(u\)，\(v\in[-1,1]\)，都有 \(|f(u)-f(v)|\leqslant |u-v|\).
\end{circlelist}
\begin{enumerate}
\item 证明：对任意的 \(x\in[-1,1]\)，\(x-1\leqslant f(x)\leqslant 1-x\)；
\item 证明：对任意的 \(u\)，\(v\in[-1,1]\)，\(|f(u)-f(v)|\leqslant 1\)；
\item 在区间 \([-1,1]\) 上是否存在满足题设条件的奇函数 \(y=f(x)\)，且使得 \(\begin{cases}
|f(u)-f(v)|<|u-v|, & u,v\in\left[0,\frac12\right],\\
|f(u)-f(v)|=|u-v|, & u,v\in\left[\frac12,1\right],
\end{cases}\) 若存在，请举一例；若不存在，请说明理由.
\end{enumerate}

\begin{solution}
\begin{enumerate}
\item 由题设 \(f(1)=0\)，对任意 \(x\in[-1,1]\)，有 \(\lvert f(x)\rvert=\lvert f(x)-f(1)\rvert\leqslant\lvert x-1\rvert=1-x\). 因此 \(-(1-x)\leq f(x)\leq1-x\)，即 \(x-1\leq f(x)\leq1-x\).
\item 不妨设 \(u\leq v\)，令 \(d=v-u\)，则 \(0\leq d\leq2\). 由题设直接得到 \(\lvert f(u)-f(v)\rvert\leq d\). 另一方面，利用 \(f(-1)=f(1)=0\) 以及 Lipschitz 条件，有 \(\lvert f(u)-f(v)\rvert\leq\lvert f(u)-f(-1)\rvert+\lvert f(1)-f(v)\rvert\leq(u+1)+(1-v)=2-d\). 所以 \(\lvert f(u)-f(v)\rvert\leq\min\{d,2-d\}\leq1\).
\item 若存在这样的奇函数，则由奇性知 \(f(0)=0\). 在第二个条件中取 \(u=\frac12\)，\(v=1\)，结合 \(f(1)=0\) 得 \(\lvert f(\frac12)\rvert=\frac12\). 但在第一个条件中取 \(u=0\)，\(v=\frac12\)，则应有 \(\lvert f(\frac12)-f(0)\rvert<\frac12\)，即 \(\lvert f(\frac12)\rvert<\frac12\)，矛盾. 因此不存在满足条件的函数.
\end{enumerate}
\end{solution}
\end{problem}
